\section{Coordinate weights and localized corrections}\label{sec:weights}

The shift will replace one scalar in each coordinate block by a scalar
from the preceding block. We first describe this operation, then choose
weights whose dependence on the input is slow and whose correction
coordinates have the support properties needed for the later tests.

For a countably infinite set $A$, let $c_0(A)$ have its supremum norm.
Write $\mathbb N=\{1,2,\ldots\}$, index levels by
$\lambda=(i,n)\in\mathbb N^2$, and put $J_\lambda=i+n$. For $j\geq1$ set
\[
 d_j=2^{-j},\qquad
 \mathcal A_j=([-1,1]\cap d_j\mathbb Z)^j,
\]
and let $\mathcal A_0$ consist of one point. The finite coordinate block at
$\lambda$ is
\[
 \Gamma_\lambda
   =\{(\lambda,j,\xi):0\leq j\leq J_\lambda,
                              \ \xi\in\mathcal A_j\}.
\]
We use
\[
 E=c_0(\mathbb N),\qquad
 U=c_0\!\left(\bigsqcup_{\lambda\in\mathbb N^2}\Gamma_\lambda\right),
 \qquad D=E\oplus_\infty U,
\]
and write $x=(t,(x^\lambda)_\lambda)$. Put $P(t,u)=u$.
Fix an enumeration of the disjoint union of the coordinate sets of $E$
and $U$. The map $\pi_j:D\to\mathbb R^j$ retains the first $j$
coordinates in this enumeration; $\pi_0$ is the unique map into
$\mathbb R^0$. Write $B=\{s\in D:\|s\|_D\leq1\}$. For $n>1$
let $\lambda^-=(i,n-1)$ and set
\[
 r_{i,1}(x)=\max\{|t_i|,\|x^{i,1}\|_\infty\},\qquad
 r_{i,n}(x)=\max\{\|x^{i,n-1}\|_\infty,\|x^{i,n}\|_\infty\}.
\]
Each $r_\lambda$ is $1$-Lipschitz. Every input coordinate occurs in at most
two of these maxima. Thus, for every fixed $x\in D$,
\begin{equation}
 r_\lambda(x)\longrightarrow0
 \quad\text{as $\lambda$ leaves finite sets}.\label{eq:radii-tail}
\end{equation}
Indeed, infinitely many radii at least $\varepsilon>0$ would require
infinitely many input coordinates of modulus at least $\varepsilon$.

The scalar replacement uses a vector $v\in\mathbb R^{\Gamma_\lambda}$
and a row $w\in(\mathbb R^{\Gamma_\lambda})^*$ with $wv=1$.
For a block vector $z\in\mathbb R^{\Gamma_\lambda}$ and an incoming
scalar $b\in\mathbb R$,
\[
 z=(z-vwz)+v(wz),\qquad
 z+v(b-wz)=(z-vwz)+vb.
\]
The first summand belongs to $\ker w$, so the operation preserves it
and replaces the scalar readout $wz$ by $b$. Along a row of blocks,
we will take $b$ from the preceding block's readout, with $t_i$
entering the first block. Section~\ref{sec:shift} constructs the inverse
for any fixed collection of such pairs. The conditions in (W1) below
combine this scalar identity with uniform bounds on the vectors and rows.

The pairs will depend on a state $s\in B$. Condition (W2) makes their
variation small after multiplication by the local radius, which bounds
the size of the input blocks involved in the replacement. This is the
estimate that will preserve the frozen shift's distance bounds when
the weights vary with the input.

The bands $j$ in each block serve a different purpose: localization.
For $j\geq1$, a correction coordinate indexed by $(\lambda,j,\xi)$
will be supported where the first $j$ state coordinates are close to
$\xi$. Whenever $j<J_\lambda$, it will also require a positive lower
bound on $r_\lambda(s)$ depending only on $j$. The first restriction
becomes sharper as $j$ grows; the second gives a fixed radial threshold
when $j$ is fixed and the level varies. We will prove these two support
bounds for the correction coordinates after constructing the weights.

\begin{lemma}[Weights at one level]\label{lem:coordinate-weights}
Let $\eta=10^{-3}$.  There are continuous maps
\[
 v^\lambda:B\to\mathbb R^{\Gamma_\lambda},\qquad
 w^\lambda:B\to(\mathbb R^{\Gamma_\lambda})^*
\]
depending on finitely many coordinates of their argument, such that
\begin{equation*}
 \|v^\lambda(s)\|_\infty\leq1,\qquad
 \|w^\lambda(s)\|_1\leq2,\qquad
 w^\lambda(s)v^\lambda(s)=1,                         \tag{W1}
\end{equation*}
and, for $s,s'\in B$,
\begin{equation*}
 \min\{r_\lambda(s),r_\lambda(s')\}
 \max\!\left\{\|v^\lambda(s)-v^\lambda(s')\|_\infty,
              \|w^\lambda(s)-w^\lambda(s')\|_1\right\}
 \leq\eta\|s-s'\|_D.                                  \tag{W2}
\end{equation*}
More precisely, the $j$th band of either weight is the product of a
coefficient $\theta_j^{J_\lambda}(r_\lambda(s))$ and a vector or row
depending only on $\pi_j s$.  If $\theta_j^{J_\lambda}(r)>0$, then
\begin{equation}
 r\leq\rho_j\quad(j\geq1),\qquad
 r\geq\ell_j\quad(j<J_\lambda),                       \label{eq:band-support}
\end{equation}
for positive constants $\rho_j,\ell_j$ defined in the proof.
\end{lemma}

\begin{proof}
Fix $j\geq1$ and $s\in B$.  For $\xi\in\mathcal A_j$ put
\[
 b^j_\xi(s)
   =\left(1-\frac{\|\pi_j s-\xi\|_\infty}{d_j}\right)_+,
 \qquad B_j(s)=\sum_{\xi\in\mathcal A_j}b^j_\xi(s).
\]
The grid includes both endpoints in every coordinate.  Every point of
$[-1,1]^j$ is within $d_j/2$ in supremum norm of a grid point, including
points on the boundary of the cube.  Thus $B_j(s)\geq1/2$; in particular,
the following row is defined everywhere on $B$:
\[
 w^j_\xi(s)=\frac{b^j_\xi(s)}{B_j(s)},\qquad
 v^j_\xi(s)
   =\min\!\left\{1,
       \left(2-\frac{\|\pi_j s-\xi\|_\infty}{d_j}\right)_+\right\}.
\]
The entries are nonnegative, $\|w^j(s)\|_1=1$ and
$\|v^j(s)\|_\infty\leq1$.  If $w^j_\xi(s)>0$, then
$\|\pi_j s-\xi\|_\infty<d_j$ and $v^j_\xi(s)=1$.
Consequently $w^j(s)v^j(s)=1$.  At band $j=0$, set both the sole
coordinate of $v^0$ and the sole entry of $w^0$ equal to $1$.

Here are explicit finite Lipschitz bounds.  Put
$N_j=|\mathcal A_j|=(2^{j+1}+1)^j$.  Since each $b^j_\xi$ and
$v^j_\xi$ has Lipschitz constant at most $d_j^{-1}$,
\[
 \|b^j(s)-b^j(s')\|_1
      \leq N_jd_j^{-1}\|s-s'\|_D.
\]
Normalizing a nonnegative vector whose $\ell_1$ norm is at least $1/2$
changes distances by at most a factor $4$ in $\ell_1$.  Indeed,
\[
 \left\|\frac b{\|b\|_1}-\frac{b'}{\|b'\|_1}\right\|_1
 \leq\frac{\|b-b'\|_1+|\|b\|_1-\|b'\|_1|}{\|b\|_1}
 \leq4\|b-b'\|_1.
\]
Thus $L_j=4N_j/d_j$ bounds both $\operatorname{Lip}(w^j)$ in
$\ell_1$ and $\operatorname{Lip}(v^j)$ in the supremum norm.

Choose $H>0$ with $\pi/H\leq\eta/2$ and inductively choose numbers
$\rho_j>0$ so that
\begin{equation}
 0<\rho_1<1,\qquad
 \rho_{j+1}<e^{-H}\rho_j,\qquad
 L_j\rho_j\leq\eta/4\quad(j\geq1).                 \label{eq:scale-input}
\end{equation}
Each constraint at stage $j$ has a positive right-hand side, so this
choice is possible.  Put $\ell_j=e^{-H}\rho_{j+1}$ for $j\geq0$.
For an integer $J\geq1$ define
$\theta^J(r)=(\theta^J_0(r),\ldots,\theta^J_J(r))$ on $[0,1]$ as
follows.  For $r\geq\rho_1$ take $\theta^J(r)=e_0$.  As $r$
decreases across $[e^{-H}\rho_j,\rho_j]$, $1\leq j\leq J$,
rotate from $e_{j-1}$ to $e_j$ by setting
\[
 \theta^J_{j-1}(r)=\cos\alpha_j(r),\qquad
 \theta^J_j(r)=\sin\alpha_j(r),\qquad
 \alpha_j(r)=\frac{\pi}{2H}\log\frac{\rho_j}{r},
\]
with all other entries zero.  Between these transition intervals keep
the most recently reached unit vector, and for
$0\leq r\leq e^{-H}\rho_J$ take $\theta^J(r)=e_J$.  The strict
inequality in \eqref{eq:scale-input} makes the intervals disjoint and ordered.
At each transition endpoint the formula agrees with the adjacent
constant value; this also proves continuity at $r=0$.  Every entry is
nonnegative,
\begin{equation}
 \sum_{j=0}^J\theta^J_j(r)^2=1,\qquad
 \sum_{j=0}^J\theta^J_j(r)\leq2.                    \label{eq:theta-sums}
\end{equation}
For $0<r\leq1$, the path $u\mapsto\theta^J(e^u)$ is
piecewise continuously differentiable and has $\ell_1$ speed at most
$(\pi/(2H))(|\sin\alpha|+|\cos\alpha|)\leq\pi/H$.
It is therefore $\pi/H$-Lipschitz in the variable $\log r$.
The schedule also gives \eqref{eq:band-support}.  For $j=0$, only
the second assertion applies.  For $j=J$, only the first assertion
applies, and $\theta_J^J(0)=1$ is allowed.

For $J=J_\lambda$ and $r=r_\lambda(s)$, concatenate the bands:
\[
 v^\lambda(s)
   =\bigl(\theta_j^J(r)v^j(s)\bigr)_{j=0}^J,
 \qquad
 w^\lambda(s)
   =\bigl(\theta_j^J(r)w^j(s)\bigr)_{j=0}^J.
\]
Every band depends on finitely many coordinates: $\pi_j s$ and
$r_\lambda(s)$ involve finite coordinate sets.  Continuity includes
the cases $r=0$ and transition endpoints.  The bounds on the band
weights and \eqref{eq:theta-sums} give
\[
 \|v^\lambda(s)\|_\infty\leq1,\quad
 \|w^\lambda(s)\|_1=\sum_{j=0}^J\theta_j^J(r)\leq2,\quad
 w^\lambda(s)v^\lambda(s)
       =\sum_{j=0}^J\theta_j^J(r)^2=1.
\]
This proves (W1).

For (W2), put $r=r_\lambda(s)$, $r'=r_\lambda(s')$,
$m=\min\{r,r'\}$ and $\delta=\|s-s'\|_D$.  If $m=0$ there is
nothing to prove.  Otherwise
\[
 m|\log r-\log r'|\leq|r-r'|\leq\delta.
\]
Change the band coefficients first, using the band weights at $s'$.
The logarithmic Lipschitz estimate and the bounds
$\|v^j(s')\|_\infty,\|w^j(s')\|_1\leq1$ show that this part of
either norm difference, multiplied by $m$, is at most
$(\pi/H)\delta\leq(\eta/2)\delta$.  Then change the band weights
with coefficients fixed at $\theta^J(r)$.  Only positive $j$ vary,
and $\theta^J_j(r)>0$ implies $m\leq r\leq\rho_j$.  Hence this
part of either norm difference, multiplied by $m$, is at most
\[
 \delta\sum_{j=1}^J\theta^J_j(r)mL_j
 \leq\frac{\eta\delta}{4}
       \sum_{j=1}^J\theta^J_j(r)
 \leq\frac{\eta\delta}{2}.
\]
The sum of the two contributions proves (W2).
\end{proof}

We record the support information in the form used by the lifting
construction.  For $s\in B$ define
\begin{equation}\label{eq:correction-coordinate}
 a_{i,0}(s)=t_i,\qquad
 a_{i,n}(s)=w^{i,n}(s)s^{i,n},\qquad
 f_{(\lambda,j,\xi)}(s)
   =\theta_j^{J_\lambda}(r_\lambda(s))v^j_\xi(s)
       \bigl(a_{i,n-1}(s)-a_{i,n}(s)\bigr).
\end{equation}
These functions will be the coordinates of the shift correction.
We first prove their support and continuity properties directly from
the weights.

\begin{lemma}[Finite dependence and local support]
\label{lem:weight-support}
For every $\gamma=(\lambda,j,\xi)$, the function $f_\gamma$ is
continuous, vanishes at $0$, and depends on a nonempty finite set
$I_\gamma$ of coordinates of $D$.  The sets can be chosen to include
the first $j$ coordinates and every coordinate used by
$r_\lambda$.  Uniformly in $\gamma$,
\begin{equation}\label{eq:correction-bounds}
 |f_\gamma(s)|\leq4r_\lambda(s)\leq4,\qquad
 \operatorname{Lip}(f_\gamma)\leq4+12\eta.
\end{equation}
Moreover
\begin{equation}
 f_\gamma(s)\ne0\quad\Longrightarrow\quad
 \begin{cases}
 \|\pi_j s-\xi\|_\infty\leq2d_j,&j\geq1,\\
 r_\lambda(s)\geq\ell_j,&j<J_\lambda,
 \end{cases}                                                   \label{eq:correction-support}
\end{equation}
where both conditions are imposed when both apply.
\end{lemma}

\begin{proof}
For $n>1$, $a_{i,n-1}$ uses only $s^{i,n-1}$, the local radius
$r_{i,n-1}(s)$, and $\pi_{J_{i,n-1}}s$.  The former radius uses
$s^{i,n-2}$ when $n>2$ and $t_i$ when $n=2$.  The current scalar
$a_{i,n}$ uses only $s^{i,n}$, $r_{i,n}(s)$, and
$\pi_{J_{i,n}}s$; for $n=1$ the predecessor scalar is $t_i$.
Since $J_{i,n-1}<J_{i,n}$, one possible common finite set for all
coordinates $\gamma$ at $\lambda=(i,n)$ is the union of the first
$J_\lambda$ enumerated coordinates, the coordinate $t_i$, and every
coordinate of the blocks $\Gamma_{i,k}$ with
$k\in\{n,n-1,n-2\}\cap\mathbb N$.  Call it $I_\gamma$.  It is
nonempty and contains the coordinates required in the statement.

The row bound (W1) gives
$|a_{i,n-1}(s)-a_{i,n}(s)|\leq4r_\lambda(s)$, also when $n=1$.
Thus $|f_\gamma(s)|\leq4r_\lambda(s)\leq4$ and
$f_\gamma(0)=0$.  To verify the uniform Lipschitz bound, fix
$s,s'\in B$ and write $\delta=\|s-s'\|_D$.  At each level
choose the pair $(\bar v^\lambda,\bar w^\lambda)$ from whichever
of $s,s'$ has the larger local radius.  If $e\in\{s,s'\}$ is the
other endpoint, (W2) gives
\[
 r_\lambda(e)\max\!\left\{
    \|v^\lambda(e)-\bar v^\lambda\|_\infty,
    \|w^\lambda(e)-\bar w^\lambda\|_1\right\}
 \leq\eta\delta;
\]
the same bound is trivial if $e$ supplied the frozen pair.  For
$e=s,s'$, replacing its weights by the frozen weights in the
correction vector
$v^\lambda(e)(a_{i,n-1}(e)-a_{i,n}(e))$ costs at most
$4\eta\delta$ from the vector and at most $2\eta\delta$
from its current and predecessor rows together.  At $n=1$ the
predecessor $t_i$ does not change.  With all weights frozen, the
correction vector is linear in the input and has norm at most
$4r_\lambda(s-s')\leq4\delta$.  Comparing the two endpoints
therefore gives $|f_\gamma(s)-f_\gamma(s')|
\leq(4+12\eta)\delta$.

If $f_\gamma(s)\ne0$, its displayed product has both
$\theta_j^{J_\lambda}(r_\lambda(s))>0$ and
$v^j_\xi(s)>0$.  For $j\geq1$, the formula for $v^j_\xi$ gives
$\|\pi_j s-\xi\|_\infty<2d_j$, which implies the first
closed inequality in \eqref{eq:correction-support}.  If $j<J_\lambda$, the second inequality
follows from \eqref{eq:band-support}.  This includes $j=0$;
no grid condition is imposed there.
\end{proof}
